\begin{corollary}[6.12.5]
\label{IV.2.s6.12.5-fr}
\textnormal{(Zariski).} --- Sea $k$ un cuerpo; para todo $k$-preesquema $X$
localmente de tipo finito sobre $k$, el conjunto de los $x\in X$ en los que
$X$ es regular \textup{(resp.} geométricamente regular sobre $k$\textup{)} es
abierto en $X$.
\end{corollary}

\begin{proof}
La afirmación del corolario relativa a la propiedad de ser regular se deduce
de (6.12.4) tomando $A=k$. La afirmación relativa a la propiedad de ser
geométricamente regular se deduce ya de (6.8.7); también se puede deducir de
(6.12.4) del modo siguiente: pongamos $k'=k^{p^{-\infty}}$ ($p$ es el
exponente característico de $k$); como el morfismo
$\operatorname{Spec}(k')\to\operatorname{Spec}(k)$ es radicial, entero y
sobreyectivo, es un homeomorfismo universal (2.4.5), por lo que el morfismo
proyección $X\otimes_kk'\to X$ es un homeomorfismo. La proyección en $X$ de
$\operatorname{Reg}(X\otimes_kk')$ es el conjunto de los puntos de $X$ en los
que $X$ es geométricamente regular sobre $k$, en virtud de (6.7.7, \emph{e)});
por tanto, ese conjunto es abierto según lo anterior.
\end{proof}
